\section{Methodology}
\label{sec:methodology}

The work follows an incremental, design-science approach. The plan front-loads
learning and builds a thin end-to-end slice (one vulnerability family, server-side
verified scoring, one attempt logged) before adding breadth, and it reuses mature
tools rather than reinventing them. Wherever machine learning appears, it is a
simple, interpretable model.

\subsection{System overview and workflow}

The platform is a standard web application. The intentionally vulnerable web
application runs as a \emph{separate}, network-isolated service, so its seeded
vulnerabilities cannot be turned against the platform, the student records, or other
students. The design keeps one clear separation: a synchronous \emph{scoring path}
that runs while the student waits (verify the exploit, then score from completion,
accuracy, and time, and return points and badges), and an asynchronous
\emph{analytics path} that runs in the background (predict who is struggling,
recommend the next challenge, and feed the instructor dashboard). Keeping the second
path off the request means scoring is never slowed and the machine-learning code can
never change a grade. Figure~\ref{fig:block} shows the deployed components.

\diagramfig{block-diagram}{Block diagram: the vulnerable service is kept physically
separate from the platform, and the ML worker runs off the request
path.}{fig:block}

\subsection{Evaluation design}

Effectiveness (O6) is assessed on three axes. For \emph{practical skill}, a pilot
cohort takes a pre-test and a post-test on unseen challenges, and the skill gain is
measured. For \emph{engagement}, a short validated questionnaire is combined with
usage metrics (attempts, time on task, and return visits). For the \emph{ML
component}, the at-risk model is validated on a held-out split of the collected
attempt data, reporting prediction accuracy and whether its recommendations match an
expert judgement of challenge difficulty. Where the cohort allows, a configuration
with the adaptive and gamified features enabled is compared against a baseline with
them disabled, following the at-risk-detection methodology of prior
work~\cite{Svabensky2024detecting}.
